\section{The weighted S-construction}\label{sec:waldhausen}

We construct the motives from complete compact flags and fixed
exact-sequence weights. No motive or K-theory functor is applied to
the input Homs.

Fix $n\geq1$ and $\mathcal A\in\nPr{(n+1)}$.

\subsection{Exact flags and their classifying modules}

Let $\operatorname{Ar}[q]$ be the poset of pairs $(i,j)$ with
$0\leq i\leq j\leq q$, ordered coordinatewise.

\begin{definition}\label{def:exact-flags}
For $H\in\nPr d$, with $d\geq1$, write
$\Fun_0(\operatorname{Ar}[q],H)$ for coherent diagrams
$x$ with specified diagonal zeros $x_{ii}=0$.
An \emph{exact flag} of length $q$ is such a
diagram together with a complete exact sequence
\[
 x_{ij}\longrightarrow x_{ik}\longrightarrow x_{jk}
       \qquad(0\leq i<j<k\leq q)
\]
above each three-vertex subdiagram. The sequences share the displayed
objects, arrows, and zero comparisons. Define their structured diagram
data by the pullback
\begin{equation}\label{eq:flag-definition-pullback}
 S_q^{\ex}(H)=
 \Fun_0(\operatorname{Ar}[q],H)
 \mathop{\times}_{\displaystyle\prod_{i<j<k}
                  \Fun_0(\operatorname{Ar}[2],H)}
             \prod_{i<j<k}\EqRow_d(H).
\end{equation}
The maps from $\EqRow_d(H)$ retain the complete coherent square and
forget the extra exactness data. The
formula specifies objects, strong maps, and their local coefficient
Homs separately.

A morphism of flags is a morphism of these complete diagrams. It preserves the complete adjunction diagrams, including the right adjoints and their compatibility with units and counits. Hence the induced Beck–Chevalley mates are equivalences. For two fixed flags
$\xi,\zeta$, let $\mathcal H_{q,H}(\xi,\zeta)\in\nPr{(d-1)}$
be their coefficient Hom. It is the small end of the coefficient
Homs of $H$ over the complete walking diagram, including the
right-adjoint arrows and their unit and counit cells.
\end{definition}

The end is formed in the coefficient category and encodes the full adjunction data, including the right-adjoint comparisons and their Beck--Chevalley compatibilities. 

By Lemma~\ref{lem:exact-row-structure}, it computes precisely the coefficient Homs of strong maps in \eqref{eq:flag-definition-pullback}, rather than merely the maps of the underlying forward diagrams.

Restriction along monotone maps gives $S_\bullet^{\ex}(H)$ a
simplicial structure. Repeated vertices give the canonical degenerate
rows. Compatible right-adjoint data for a fixed left adjoint form a
contractible space by the matrix-duality construction of
Section~\ref{sec:weights}. Hence restriction and composition of flags
give coherent operators. In particular,
$S_0^{\ex}(H)$ is the zero flag and $S_1^{\ex}(H)=H$ as diagram data.

Let $\mathcal F_q(\mathcal A)$ be a small family representing all
complete flags with compact interval objects, together with their
strong maps and coefficient Homs.
For $n\geq2$, form the coefficient matrix
\[
 \mathcal R_{n,q}(\xi,\zeta)
   =\mathsf P_{n-1}
      \bigl(\mathcal H_{q,\mathcal A}(\xi,\zeta)\bigr)
       \in\nPr{(n-1)}.
\]
Composition is induced by the lower free external product and
composition of the ends. Define its raw module by
\begin{equation}\label{eq:waldhausen-input-flags}
 Q_{n,q}(\mathcal A)=
 \begin{cases}
 \Fun(\mathcal F_q(\mathcal A)^{\op},\Sp),&n=1,\\
 \Pmod_{\nPr{(n-1)}}(\mathcal R_{n,q}),&n\geq2.
 \end{cases}
\end{equation}
At $n=1$ the argument of $\Fun$ is the ordinary category of complete
compact flags. At every index the outer represented values are
arbitrary, while local values satisfy the compact-image condition.
The free module property gives its universal property after every
coefficient parameter. In particular $Q_{n,1}(\mathcal A)=
\mathsf P_n(\mathcal A)$. This is the coefficient-matrix variant
of Definition~\ref{def:raw-presheaves}. Restriction gives simplicial
operators preserving represented compact generators.

\begin{definition}\label{def:flag-weights}
Let
\[
 G_{n,q}=\mathsf F_n(\operatorname{Ar}[q])
       \amalg_{\one_n^{\oplus(q+1)}}0,
 \qquad B_n=G_{n,2}.
\]
The map from $\one_n^{\oplus(q+1)}$ selects the diagonal objects.
The universal exact sequence defines $B_n\to W_n$, where $W_n$ is the
module of Proposition~\ref{def:presentable-exact-classifier}. For
$\mathcal T_q=\{(i,j,k)\mid0\leq i<j<k\leq q\}$, set
\begin{equation}\label{eq:flag-weight-pushout}
 T_{n,q}=
 G_{n,q}
 \amalg_{\displaystyle\coprod_{\tau\in\mathcal T_q}B_n}
             \displaystyle\coprod_{\tau\in\mathcal T_q}W_n.
\end{equation}
The other map from each copy of $B_n$ selects the corresponding
three-vertex subdiagram. These are pushouts of coefficient modules.
\end{definition}

\begin{proposition}\label{prop:flag-weight-up}
There are natural parameter-compatible equivalences
\begin{equation}\label{eq:flag-weight-up}
 \Fun^L_{\nPr{(n-1)}}(T_{n,q},E)
       \simeq S_q^{\ex}(E).
\end{equation}
The weights form a cosimplicial module $T_{n,\bullet}$ whose maps
preserve compact objects. Moreover,
\[
 T_{n,0}=0,\qquad T_{n,1}\simeq\one_n,\qquad
 T_{n,2}\simeq W_n.
\]
\end{proposition}


\begin{proof}
The unrestricted exact-sequence classifier
of Corollary~\ref{cor:presentable-exact-classifier}
and the free module property identify
\eqref{eq:flag-weight-pushout} with
\eqref{eq:flag-definition-pullback},
including coefficient parameters.

Restriction of flags induces the cosimplicial
structure. Its maps preserve compact objects
by \eqref{eq:compact-exact-classifier},
since the universal interval objects are compact.
The low-degree identifications follow directly
from the defining pushout.
\end{proof}


At $n=1$, $T_{1,q}$ is the Ind-completion of the perfect category of
a linearly oriented $q$-vertex spectral quiver. Indeed, the usual
gap-diagram equivalence
$S_q^{\ex}(E)\simeq\Fun([q-1],E)$ identifies their representing
properties \citep{Waldhausen}.

\subsection{Recognition from degree one}

The following arguments take place in ordinary presentable categories
and their coefficient Hom objects. They also apply in the parameter
modules of Lemma~\ref{lem:parameter-modules}.


\begin{lemma}\label{lem:exact-row-cofiber}
Let $n\geq 1$ and $E\in\nPr{n}$.
Every internal exact sequence
\(a\xrightarrow{i}b\xrightarrow{p}c\)
in $E$ determines a coherent cofiber square in the underlying presentable category of $E$.
This quotient property is compatible with all strong
parameter diagrams.
\end{lemma}

\begin{proof}
For $n=1$ this is the definition. For $n\geq2$, precomposition by
$p$ is fully faithful, with inverse on its image given by
precomposition by $j$. For $f:b\to z$, postcomposition of the central
exact sequence gives
\[
 fir\longrightarrow f\longrightarrow fjp.
\]
If $fi=0$, the first term is zero.
By induction, the second arrow is an
equivalence.
Thus \(f\) naturally factors through \(p\).
\end{proof}

\begin{lemma}\label{lem:flag-first-row}
The first-row chain
\[
 x_{01}\longrightarrow x_{02}\longrightarrow\cdots
       \longrightarrow x_{0q}
\]
of an exact flag in $E$ determines its interval
quotients and coherent comparisons up to contractible choice. This
holds with every coefficient parameter. A specified map of first-row
chains has either no lift to a morphism of complete flags or a
contractible space of lifts.
\end{lemma}

\begin{proof}
Each $x_{ij}$ is the quotient of $x_{0i}\to x_{0j}$ by
Lemma~\ref{lem:exact-row-cofiber}. The pushout universal property
determines the quotient objects, arrows, and zero comparisons
together. The compatible right-adjoint data are contractible when
they exist.
A map of first-row chains determines a map of the quotient diagrams.
It lifts to a morphism of complete flags exactly when its canonical
right-adjoint mates are equivalences. This is a property of those
comparisons. If it holds, the compatible adjunction data give the
contractible space of lifts. The construction in each parameter
module gives the asserted parameter compatibility.
\end{proof}


For the flag matrices, $\operatorname{Fun}_n
(\mathcal F_p(\mathcal A),S_q^{\ex}(E))$ means assignments of target
flags with coherent unit and composition comparisons and the local
compact-image condition. For $n\geq2$ its local maps use
$\mathsf P_{n-1}(\mathcal H_{p,\mathcal A}(\xi,\zeta))$ and the
target flag coefficient Homs. At $n=1$ these are ordinary functors.
Write
$\operatorname{Fun}_{n,\mathrm{simp}}
(\mathcal F_\bullet(\mathcal A),S_\bullet^{\ex}(E))$
for the end imposing compatibility with simplicial operators,
taken on objects and coefficient Homs.

\begin{theorem}\label{thm:flag-recognition}
Restriction to degree one induces a parameter-compatible equivalence
\begin{equation}\label{eq:flag-recognition}
 \operatorname{Fun}_{n,\mathrm{simp}}
    (\mathcal F_\bullet(\mathcal A),S_\bullet^{\ex}(E))
       \simeq\Add_n^\omega(\mathcal A,E).
\end{equation}
\end{theorem}


\begin{proof}
An additive diagram $F$ preserves exact flags,
including their adjunction data and coherent
comparisons. It therefore induces a simplicial
diagram $S_\bullet^{\ex}(F)$.
Conversely, let $H_\bullet$ be a simplicial
diagram as in the statement, and put $F=H_1$.
Degree zero gives $F(0)=0$.
For an exact sequence $e=(a\to b\to c)$,
the faces identify the vertices of $H_2(e)$
with $F(a),F(b),F(c)$.

By Lemma~\ref{lem:canonical-row-maps}
and \eqref{eq:row-evaluation-adjunctions},
its forward arrows and zero comparison
agree with those of $F(e)$.
Uniqueness of adjunction data
then identifies the remaining structure.
Thus $F$ is additive and
$H_2\simeq S_2^{\ex}(F)$.

For $q\geq3$, the interval faces identify
the first-row chain of $H_q$ with the image
under $F$ of the source first-row chain.
The resulting comparisons are equivalences,
so their right-adjoint mates are invertible.
Lemma~\ref{lem:flag-first-row} therefore
extends them uniquely up to contractible choice
to equivalences $H_q\simeq S_q^{\ex}(F)$.
This uniqueness ensures compatibility with
all simplicial operators.

The same argument in coefficient diagram modules
identifies the coefficient Homs and proves
parameter compatibility.
\end{proof}


\subsection{The weighted realization and its universal property}

\begin{definition}[The weighted Waldhausen \(S\)-construction]\label{def:weighted-s-construction}
With $Q_{n,q}(\mathcal A)$ as in \eqref{eq:waldhausen-input-flags}, define
\begin{equation}\label{eq:weighted-s-construction}
 \operatorname{Wald}_n(\mathcal A)
       =\int^{[q]\in\Delta}T_{n,q}\otimes Q_{n,q}(\mathcal A).
\end{equation}
This coend is the realization of the two-sided simplicial bar
\begin{equation}\label{eq:weighted-s-bar}
 \mathcal B_p=
 \coprod_{[q_0]\to\cdots\to[q_p]\ \mathrm{in}\ \Delta}
       T_{n,q_0}\otimes Q_{n,q_p}(\mathcal A).
\end{equation}
The first face uses the cosimplicial action on $T_{n,\bullet}$.
The last face uses restriction of input flags. The other faces
compose monotone maps, and the degeneracies insert identities.
\end{definition}

The structure maps preserve represented compact generators and
compact coefficient actions, so the tensor products and module
colimit belong to $\nPr n$. Their free properties allow all
coefficient-linear left adjoints in the mapping-out property.
Since $T_{n,1}=\one_n$ and $Q_{n,1}=\mathsf P_n(\mathcal A)$,
the degree-one map gives a
canonical diagram
\[
 U^\omega_{\mathcal A}\in
 \operatorname{Fun}_n^\omega(\mathcal A,
                    \operatorname{Wald}_n(\mathcal A)).
\]

\begin{theorem}\label{thm:weighted-s-universal}
The diagram $U^\omega_{\mathcal A}$ extends to a filtered-continuous additive
diagram on all objects of $\mathcal A$. For every $E\in\nPr n$,
restriction induces a parameter-compatible equivalence
\begin{equation}\label{eq:weighted-s-universal}
 \Fun^L_{\nPr{(n-1)}}(\operatorname{Wald}_n(\mathcal A),E)
       \simeq\Add_{n}^\omega(\mathcal A,E)\cong\Add_{n,\mathrm{filt}}(\mathcal A,E).
\end{equation}
Consequently there is a canonical equivalence
$\operatorname{Wald}_n(\mathcal A)\simeq\M_n(\mathcal A)$, natural
for module morphisms and their coefficient parameter diagrams.
\end{theorem}

\begin{proof}
Let $E^{\mathcal R_{n,p}}$ be the module of diagrams for the raw
flag coefficient matrix, as in Lemma~\ref{lem:parameter-modules}.
At $n=1$ use $\Fun(\mathcal F_p(\mathcal A),E)$.
Free evaluation, tensor currying, and
Proposition~\ref{prop:flag-weight-up} give
\[
\begin{aligned}
 \Fun^L_{\nPr{(n-1)}}
    (T_{n,q}\otimes Q_{n,p}(\mathcal A),E)&\simeq\Fun^L_{\nPr{(n-1)}}(T_{n,q},E^{\mathcal R_{n,p}})\\
 &\simeq S_q^{\ex}(E^{\mathcal R_{n,p}})\\
 &\simeq\operatorname{Fun}_n
       (\mathcal F_p(\mathcal A),S_q^{\ex}(E)).
\end{aligned}
\]
The last step curries complete rows and their right-adjoint data
pointwise in the diagram module.
Mapping out of the coend gives the end
\[
 \Fun^L_{\nPr{(n-1)}}(\operatorname{Wald}_n(\mathcal A),E)
 \simeq\operatorname{Fun}_{n,\mathrm{simp}}
      (\mathcal F_\bullet(\mathcal A),S_\bullet^{\ex}(E)).
\]
Theorem~\ref{thm:flag-recognition} identifies the right side with
$\Add_n^\omega(\mathcal A,E)$.
\end{proof}
